\documentclass[border=16pt]{standalone}
\usepackage{ctex}
\usepackage[european, straightvoltages]{circuitikz}
\begin{document}
\begin{circuitikz}[line width=0.8pt, font=\normalsize]

% fig-ge-08 共射放大电路：直流通路与微变等效电路工程图

% ==================== 顶部两栏标题 ====================
\node[anchor=south, font=\large\bfseries] at (0.3, 6.7) {直流通路（电容视为断开）};
\node[anchor=south, font=\large\bfseries] at (8.0, 6.7) {微变等效（交流通路）};

% ==================== 中间虚线分界与两侧竖排文字 ====================
\draw[dashed, gray!70, line width=0.8pt] (4.0, 0.4) -- (4.0, 6.2);
\node[rotate=90, font=\small, text=gray!25!black] at (3.64, 3.3) {大写算工作点};
\node[rotate=90, font=\small, text=gray!25!black] at (4.36, 3.3) {小写算放大};

% ==================== 左栏：直流通路 ====================
% 三极管放置
\node[npn] (Q) at (0.6, 2.5) {};

% 发射极垂直引下接地
\draw (Q.E) -- ++(0,-0.7) node[ground]{};

% 集电极引线：垂直引出 T 接点，直连 Rc 和电源
\draw (Q.C) -- ++(0, 0.6) coordinate (nodeC);
\draw (nodeC) to[R, l_=$R_c$] ++(0, 1.6) coordinate (topVcc);
\draw (topVcc) -- ++(0, 0.8) node[above, font=\normalsize] {$V_{CC}$};
\node[circ] at (topVcc) {};

% Rb 支路：自电源节点水平向左折向 Rb
\coordinate (topRb) at (-1.6, 4.7);
\draw (topVcc) -| (topRb) to[R, l_=$R_b$] (-1.6, 3.1);

% 基极引线：自 Q.B 直角折线连接到 Rb 下端
\draw (Q.B) -| (-1.6, 3.1);

% 集电极 T 形分支实心接点
\node[circ] at (nodeC) {};

% 输出耦合电容断开支路：自集电极节点向右水平引出
% 使用 circuitikz 现成 C 元件，两极板中间断开
\draw (nodeC) -- ++(0.9, 0) coordinate (capStart) to[C] ++(1.1, 0) coordinate (capEnd) -- ++(0.5, 0);
\coordinate (capMid) at ($(capStart)!0.5!(capEnd)$);
\node[above=8pt, font=\small] at (capMid) {$C$};
\node[below=8pt, font=\small] at (capMid) {（断开）};

% Q 点标注（大写）：居中在左栏电路下方
\node[anchor=north, font=\small] at (0.3, 0.7)
  {Q 点（大写）：$I_{BQ}$、$I_{CQ}$、$U_{CEQ}$};


% ==================== 右栏：微变等效电路 ====================
% 高度与左栏对齐：母线高 4.7，地线高 2.1
% 向右平移 0.4，为中间虚线与竖排文字留出充裕对称空间
\coordinate (inB)    at (5.1, 4.7);
\coordinate (inE)    at (5.1, 2.1);
\coordinate (rbeTop) at (5.8, 4.7);
\coordinate (rbeBot) at (5.8, 2.1);

% 输入端口 B 与 E
\draw (rbeTop) -- (inB) node[left, font=\small] {$B$};
\node[circ] at (inB) {};
\draw (rbeBot) -- (inE) node[left, font=\small] {$E$};
\node[circ] at (inE) {};

% 输入基极电流 Ib（标注在导线上方）
\draw[-{Stealth[length=4pt, width=3pt]}, blue!80!black] (5.2, 4.95) -- (5.7, 4.95) node[midway, above, font=\scriptsize] {$I_b$};

% r_be 支路：r_be 接在输入与地之间
\draw (rbeTop) to[R, l_=$r_{be}$] (rbeBot);
\node[circ] at (rbeTop) {};
\node[circ] at (rbeBot) {};
\draw (rbeBot) node[ground]{};

% 输出侧坐标
\coordinate (srcTop) at (7.4, 4.7);
\coordinate (srcBot) at (7.4, 2.1);
\coordinate (roTop)  at (9.1, 4.7);
\coordinate (roBot)  at (9.1, 2.1);
\coordinate (outC)   at (9.7, 4.7);
\coordinate (outE)   at (9.7, 2.1);

% 顶部集电极输出母线
\draw (srcTop) -- (roTop) -- (outC) node[right, font=\small] {$C$};
\node[circ] at (outC) {};

% 底部发射极公共地线
\draw (srcBot) -- (roBot) -- (outE) node[right, font=\small] {$E$};
\node[circ] at (outE) {};

% 受控电流源 βIb：垂直向下，标签置于右侧且留有充足间距
\draw (srcTop) to[cI, l=$\beta I_b$] (srcBot);
\node[circ] at (srcTop) {};
\node[circ] at (srcBot) {};

% ro 支路：并联在输出侧
\draw (roTop) to[R, l=$r_o$] (roBot);
\node[circ] at (roTop) {};
\node[circ] at (roBot) {};

% 输出侧接地（直流源按交流接地，发射极为参考地）
\draw (roBot) node[ground]{};

% 右侧参数列表说明
\node[anchor=west, font=\small, align=left] at (10.4, 3.4)
  {$A_u$ 电压放大\\[8pt]
   $r_i$ 输入电阻\\[8pt]
   $r_o$ 输出电阻\\[8pt]
   直流源按交流接地};

\end{circuitikz}
\end{document}
